% answers_short.tex —— 2015 年·浙江卷·文科 · 简短答案速查
% 由 tools/make_short_answers.py 自 answers.tex 的 \gksolhead/\gkanswer 提取，
% 并与 2015/2015_数学真题_浙江文_答案要点.md 逐题比对后生成，勿手改（改源 MD 后重跑）。

\begin{gknotice}
\textbf{使用说明}：本节只给最终答案便于快速核对，完整推导见解析版。答案为编写者独立推导，\textbf{非官方标准答案}。
\end{gknotice}

\gkpart{选择题}
\begin{center}\begin{tabular}{ccccccccc}
\toprule
题号 & 1 & 2 & 3 & 4 & 5 & 6 & 7 & 8 \\
\midrule
答案 & A & C & D & A & D & B & C & B \\
\bottomrule
\end{tabular}\end{center}

\gkpart{填空题}
\begin{itemize}[leftmargin=2.6em]
  \item[\textbf{9.}] \(-\dfrac12\)；\(3\sqrt3\)
  \item[\textbf{10.}] \(a_1=\dfrac23\)，\(d=-1\)
  \item[\textbf{11.}] \(\pi\)；\(\dfrac{3-\sqrt2}{2}\)
  \item[\textbf{12.}] \(-\dfrac12\)；\(2\sqrt6-6\)
  \item[\textbf{13.}] \(\dfrac{2\sqrt3}{3}\)
  \item[\textbf{14.}] \(15\)
  \item[\textbf{15.}] \(\dfrac{\sqrt2}{2}\)
\end{itemize}

\gkpart{解答题}
\begin{itemize}[leftmargin=2.6em]
  \item[\textbf{16.}] (1) \(\dfrac25\)；(2) \(9\)
  \item[\textbf{17.}] (1) \(a_n=2^n\)，\(b_n=n\)；(2) \(T_n=(n-1)2^{n+1}+2\)
  \item[\textbf{18.}] (1) 成立（建系得 \(\overrightarrow{A_1D}\) 与平面 \(A_1BC\) 的法向量平行）；(2) \(\dfrac{\sqrt7}{8}\)
  \item[\textbf{19.}] (1) \(A(2t,t^2)\)，\(B\left(\dfrac{2t}{1+t^2},\dfrac{2t^2}{1+t^2}\right)\)；(2) \(S=\dfrac{t^3}{2}\)
  \item[\textbf{20.}] (1) $g(a)=\begin{cases}\dfrac{a^2}{4}+a+2,&a\leqslant-2,\\1,&-2<a<2,\\\dfrac{a^2}{4}-a+2,&a\geqslant2.\end{cases}$；(2) \(b\in[-3,\,9-4\sqrt5]\)
\end{itemize}
